\clearpage
\section{Metodyka badań}
\label{chap:metodyka}

Metodyka obejmuje porównanie kompletnych procedur predykcyjnych oraz dwa badania kontrolowane: ważenia klas i~regularyzacji MLP. Wspólny podział zewnętrzny wyznacza populację oceny, natomiast walidacja wewnętrzna służy selekcji konfiguracji i~łączeniu predykcji.

\subsection{Projekt badania i jednostki analizy}
\label{sec:projekt-badania}

Badanie główne obejmuje siedem procedur neuronowych oraz cztery odniesienia konwencjonalne opisane w~rozdziale~\ref{chap:metody}. Każda para zbiór--metoda podlega ocenie na wspólnym zewnętrznym zbiorze testowym danego zadania. Porównaniu podlegają kompletne procedury, obejmujące przygotowanie cech, uczenie, wybór hiperparametrów i~łączenie predykcji. Różnica między dwiema takimi procedurami nie identyfikuje wpływu pojedynczego składnika architektury.

Jednostką porównania między zadaniami jest zbiór danych. Foldy wewnętrzne, konfiguracje hiperparametrów i~warunki raportowania są zależnymi elementami procedury ocenianej na tym samym zadaniu; nie zwiększają liczby niezależnych jednostek~\cite{demsar2006comparisons}. W~badaniach mechanizmów jednostką dopasowania porównań jest natomiast para warunków o~wspólnym podziale, podzbiorze i~ziarnie. Rozdzielenie tych poziomów określa zakres agregacji i~interpretację niepewności w~sekcji~\ref{sec:analiza-plan}.

Dwa badania uzupełniające wprowadzają kontrolowane zmiany przy zachowaniu pozostałych ustawień. Badanie~A dotyczy ważenia klas, a~badanie~B --- regularyzacji MLP i~liczby przykładów treningowych. Interpretacja tych porównań pozostaje odrębna od rankingu badania głównego. Ich wyniki testowe nie stanowią podstawy zmiany porównywanych procedur.

\subsection{Dobór i audyt zbiorów danych}
\label{sec:dane-audyt}

Dobór dwunastu zadań jest celowy. Ma uwzględnić różne liczności, typy cech, liczbę klas oraz relację liczby cech do liczby przykładów. Tabela~\ref{tab:rejestr-zadan} przedstawia deklarowane źródła i~role danych. Wskazane własności uzasadniają dobór, lecz nie są zmierzoną oceną trudności ani dowodem niezależności źródeł. HIGGS oznacza wskazany podzbiór OpenML, a~nie pełny pierwotny zbiór. Wersja~1 przypisana konfiguracji źródła UCI jest oznaczeniem lokalnej definicji; nie gwarantuje niezmienności pobieranego archiwum.

\begin{table}[tbp]
\centering\small
\caption{Źródła i~role przyjętych zadań. Liczba klas wynika z~definicji zadania.}
\label{tab:rejestr-zadan}
\begin{tabularx}{\linewidth}{P{0.23\linewidth}P{0.23\linewidth}cY}
\toprule
Zbiór & Źródło i~wersja & Klasy & Rola w~projekcie \\
\midrule
Adult & OpenML 1590, v2 & 2 & Cechy mieszane \\
Covertype & OpenML 1596, v4 & 7 & Duża liczność, wiele klas \\
Credit Card Fraud & OpenML 1597, v1 & 2 & Rzadkie zdarzenie \\
HIGGS, podzbiór & OpenML 23512, v2 & 2 & Cechy numeryczne \\
Amazon Employee Access & OpenML 4135, v1 & 2 & Kategorie, rzadka decyzja \\
Santander Customer Satisfaction & Oficjalny lokalny plik treningowy & 2 & Niezrównoważenie klas \\
Online Shoppers & UCI 468 & 2 & Cechy mieszane \\
Diabetes 130-US & UCI 296 & 3 & Powtarzające się pobyty pacjentów \\
Helena & OpenML 41169, v1 & 100 & Wiele klas, ograniczona semantyka cech \\
German Credit & UCI 144 & 2 & Mała próba, cechy mieszane \\
Ionosphere & UCI 52 & 2 & Mała próba, cechy numeryczne \\
Bioresponse & OpenML 4134, v1 & 2 & Dużo cech względem liczby przykładów \\
\bottomrule
\end{tabularx}
\end{table}

Definicja zadania obejmuje znaczenie etykiety, populację predykcji, mapowanie klas, typy kolumn i~sposób oznaczania braków. Ocena poprawności danych uwzględnia stałe cechy, duplikaty, identyfikatory oraz dostępność predyktorów w~momencie podejmowania decyzji. Reguły wyłączeń są niezależne od wyników porównywanych modeli; zmiana populacji na podstawie jakości testowej wprowadzałaby dodatkowy poziom selekcji.

W~Diabetes identyfikator pacjenta służy grupowaniu, a~identyfikatory pacjenta i~pobytu nie są predyktorami. W~Santander identyfikator wiersza jest odrębny od cech. Dla Credit Card Fraud istotne są interpretacja kolumny \texttt{Time} i~ograniczenia losowego podziału, dla Amazon --- powtarzające się encje, a~dla Covertype --- możliwe zależności przestrzenne. Anonimowość cech Heleny ogranicza ocenę ich znaczenia i~dostępności. Sama zgodność schematu pliku nie rozstrzyga tych kwestii.

\subsection{Podziały i przepływ informacji}
\label{sec:podzialy}

W~badaniu głównym wykorzystuje się pierwszy podział stratyfikowanej procedury trójfoldowej z~tasowaniem i~ziarnem~42. Około dwie trzecie obserwacji przypada na część przeznaczoną na uczenie i~walidację, dalej nazywaną częścią rozwojową, a~jedna trzecia na test. Pozostałe dwa podziały zewnętrzne nie są wykorzystywane jako powtórzenia eksperymentu. W~części rozwojowej zwykłe procedury uczone wykorzystują osiem foldów wewnętrznych w~jednym powtórzeniu. Automatyczne przechodzenie do powtarzanej walidacji dla małych danych zostało wyłączone.

Predykcje OOF (ang. out-of-fold) są wyznaczane dla obserwacji spoza części treningowej danego foldu. W~lokalnych MLP i~FT-Transformer etykiety walidacyjne nie uczestniczą w~obliczaniu gradientów, ale służą wyborowi epoki i~wcześniejszemu zatrzymaniu uczenia. OOF nie jest więc oceną na danych całkowicie niewykorzystanych przy dopasowaniu procedury. Predykcje te służą do wyboru konfiguracji i~wag zespołu. Zewnętrzny zbiór testowy służy niezależnej ocenie wybranej procedury~\cite{cawley2010selection}.

Dla Diabetes stosuje się stratyfikowane podziały grupowe na obu poziomach. Wszystkie pobyty jednego pacjenta muszą należeć do tej samej strony danego podziału. Kryteria poprawności podziału obejmują rozłączność grup, obecność klas oraz dokładnie jednokrotne przypisanie obserwacji do walidacji wewnętrznej. Udziały i~proporcje klas przy podziale grupowym mogą różnić się od idealnie równych; raport obejmuje rzeczywiste liczności.

Schemat~\ref{fig:podzial-selekcja} przedstawia dopuszczalny przepływ informacji. Jednokrotny podział ogranicza koszt, ale nie pozwala oszacować zmienności wynikającej ze zmiany podziału. Brak procedury czasowej oznacza również, że wyniki nie stanowią dowodu odporności na zmianę rozkładu w~czasie.

\plannedfigure
{figures/ch04/podzial-i-selekcja.pdf}
{Schemat podziału danych, wyboru konfiguracji i~końcowej oceny; opracowanie własne.}
{fig:podzial-selekcja}
{\textbf{Część rozwojowa}\quad (około $2/3$ obserwacji)\\
Osiem foldów $\longrightarrow$ predykcje OOF\\
$\longrightarrow$ wybór konfiguracji i~wag zespołu\\[0.8em]
\textbf{Część testowa}\quad (około $1/3$ obserwacji)\\
Predykcje wybranej procedury $\longrightarrow$ ocena jakości\\[0.8em]
Etykiety testowe nie uczestniczą w~selekcji.\\
IID KNN: LOO.\quad Stały TabPFN: pełny kontekst rozwojowy, bez OOF.}

\subsection{Przygotowanie cech i wyjątki proceduralne}
\label{sec:przygotowanie-procedura}

Wspólne dane wejściowe zachowują wartości i~typy cech oraz ustalone mapowanie klas. Nie narzuca się wszystkim metodom identycznego kodowania ani skalowania. Przygotowanie danych stanowi część badanego potoku, a~wyuczone przekształcenia zwykłych modeli są dopasowywane wyłącznie na treningu odpowiedniego foldu. Obsługa braków i~nieznanych kategorii musi być określona przed transformacją walidacji i~testu. Odmienne przetwarzanie kategorii jest właściwością porównywanych procedur, a~nie wyizolowanym czynnikiem eksperymentalnym.

Wyjątek stanowi natywny KNN w~zadaniach traktowanych jako IID. Korzysta on z~predykcji leave-one-out, dalej LOO, pojedynczego modelu potomnego. Etykieta przewidywanej obserwacji zostaje wyłączona z~jej sąsiedztwa, ale przygotowanie cech obejmuje całą rozwojową populację wejściową. Nie jest to ponowne dopasowanie wszystkich przekształceń po usunięciu każdego wiersza. Kategorie KNN kodowane są ustalonym kodowaniem one-hot; efektywna liczba sąsiadów i~ewentualne ograniczenie banku przez budżet należą do metadanych wykonania. W~Diabetes KNN korzysta z~ośmiu foldów grupowych.

Stały TabPFN-3.5 wykorzystuje pełną część rozwojową jako oznakowany kontekst. Nie tworzy OOF, nie przechodzi strojenia i~nie otrzymuje zespołu foldów. Jego osiem estymatorów inferencji nie jest ośmioma modelami walidacji. Porównanie obejmuje lokalny model o~ustalonej wersji wag; wiedza uzyskana podczas pretrenowania pozostaje zewnętrznym zasobem tej procedury.

\subsection{Strojenie, losowość i budżety}
\label{sec:strojenie}

Dla każdej z~dziesięciu metod podlegających strojeniu przyjęto konfigurację domyślną i~25 losowych kandydatów. Przestrzenie sześciu lokalnych integracji pochodzą z~plików YAML; KNN, RealMLP, ModernNCA i~TabM korzystają z~generatorów przypiętej biblioteki. Zapisywany jest efektywny zestaw parametrów po zastosowaniu generatorów i~nadpisań. Rozdzielenie pochodzenia przestrzeni jest istotne dla odtworzenia rzeczywiście ocenionych kandydatów.

Ziarno wyszukiwania wynosi~42. Ziarno zewnętrzne wpływa także na inicjalizację modeli, natomiast ziarna wewnętrznych podziałów są określane przez bibliotekę i~zapisywane w~artefaktach. Procedura nie korzysta z~jednego wspólnego ziarna dla wszystkich operacji. TabPFN ma stałe ziarno inferencji~0. Ziarna badań mechanizmów stanowią odrębny element definicji porównania sparowanego.

Budżet pojedynczej konfiguracji wynosi 3600~s, łącznie dla jej foldów i~przygotowania danych. Jest to limit przekazywany bibliotece, a~nie twardy termin zakończenia procesu. W~sześciu uczonych metodach neuronowych ograniczenie czasu stanowi dodatkowe kryterium zakończenia uczenia, obok maksymalnej liczby epok i~reguł opartych na walidacji. Procedura zatrzymuje uczenie z~zachowaniem stanu modelu dostępnego do predykcji; samo zatrzymanie budżetowe nie oznacza niepowodzenia dopasowania. Kontrola czasu odbywa się wewnątrz procedur uczenia, dlatego rzeczywisty czas może przekroczyć przydział. Porównanie dotyczy jakości osiąganej przy wspólnym nominalnym budżecie, a~nie po jednakowej liczbie aktualizacji parametrów. Jednakowa liczba prób nie zapewnia jednakowego czasu ani równoważnej pojemności przestrzeni.

\subsection{Warunki raportowania D, T i T+E}
\label{sec:warunki-raportowania}

Dla zwykłej konfiguracji predykcję testową stanowi średnia prawdopodobieństw z~ośmiu modeli foldów. Nie następuje dodatkowe dopasowanie wybranej konfiguracji na całej części rozwojowej. Ten sposób wykorzystania zapisanych predykcji jest inspirowany organizacją ocen w~TabArena~\cite{erickson2025tabarena}; niniejsze badanie ma jednak własny zakres i~jawne wyjątki.

\begin{table}[tbp]
\centering\small
\caption{Warunki raportowania oraz konsekwencje wyjątków. Wszystkie decyzje o~wyborze wykorzystują część rozwojową.}
\label{tab:warunki}
\begin{tabularx}{\linewidth}{P{0.14\linewidth}YY}
\toprule
Warunek & Wybór & Predykcja testowa \\
\midrule
D & Konfiguracja domyślna & Średnia foldów; dla IID KNN predykcja natywnego modelu \\
T & Najlepsza jakość OOF spośród udanych kandydatów, włącznie z~domyślnym & Zapisana predykcja wybranej konfiguracji \\
T+E & Zachłanny dobór wag na OOF, budżet do~100 kroków, ziarno~42 & Ważona suma zapisanych prawdopodobieństw konfiguracji \\
Stały TabPFN & Brak wyboru konfiguracji i~wag & Predykcja z~pełnym kontekstem rozwojowym; wyłącznie D \\
\bottomrule
\end{tabularx}
\end{table}

Dla konfiguracji $c$ podlegającej ośmiofoldowej walidacji niech $p_{c,k}(x)$ oznacza predykcję modelu foldu~$k$. Warunki D i~T używają $\bar p_c(x)=\frac{1}{8}\sum_{k=1}^{8}p_{c,k}(x)$, odpowiednio dla konfiguracji domyślnej i~wybranej. Warunek T+E ma postać
\begin{equation}
p_{\mathrm{T+E}}(x)=\sum_{c\in\mathcal{C}}\omega_c\bar p_c(x),
\qquad \omega_c\geq0,\quad \sum_{c\in\mathcal{C}}\omega_c=1,
\label{eq:zespol-konfiguracji}
\end{equation}
gdzie $\mathcal{C}$ jest zbiorem dopuszczonych konfiguracji. Wagi $\omega_c$ są wybierane na OOF. Łączenie konfiguracji dodaje więc poziom agregacji do już istniejącego uśredniania foldów; zespół wewnętrzny TabM pozostaje częścią $p_{c,k}$. Dla natywnego IID KNN w~miejsce $\bar p_c$ używana jest predykcja pojedynczego dopasowania, a~stały TabPFN występuje tylko w~D.

W~przypiętej implementacji zachowywany jest najkrótszy prefiks trajektorii zachłannego wyboru osiągający najlepszą ocenę OOF. Budżet 100 kroków nie oznacza 100 zachowanych wyborów ani 100 różnych konfiguracji; powtórny wybór tego samego kandydata zwiększa jego udział. Różnica T+E względem T mierzy korzyść z~tej dodatkowej procedury selekcji i~agregacji, przy wspólnym koszcie znalezienia kandydatów.

Wybór T i~wag T+E następuje przed wykorzystaniem etykiet testowych do obliczenia metryk. Wynik testowy nie służy wyborowi warunku raportowania. Brak kompletu wymaganych dopasowań i~predykcji lub awaria konfiguracji domyślnej uniemożliwiają pełny raport. Konfiguracja, której uczenie zakończyło się poprawnie wskutek ograniczenia czasu, pozostaje dopuszczona do raportowania oraz wyboru T i~T+E, o~ile spełnia te same wymagania kompletności co pozostali kandydaci. Jawnie nieudane konfiguracje losowe mogą zostać pominięte w~wyborze dopiero po zakończeniu prób; ich liczba i~koszt pozostają częścią opisu kompletności.

\subsection{Metryki i reguły decyzji}
\label{sec:metryki}

Dla dziewięciu zadań binarnych główną metryką wyboru jest ROC AUC, dla której większa wartość oznacza lepsze uporządkowanie. Dla trzech zadań wieloklasowych jest to strata logarytmiczna, dla której mniejsza wartość jest korzystniejsza:
\begin{equation}
L_{\log}=-\frac{1}{n}\sum_{i=1}^{n}\log p_{i,y_i}.
\label{eq:logloss-badanie}
\end{equation}
Wartość dotyczy prawdopodobieństwa przypisanego właściwej klasie. Poprawne obliczenie miary wymaga zgodności kolejności kolumn predykcji z~mapowaniem klas; sama zgodność wymiarów nie wystarcza.

Raport uzupełniają accuracy, balanced accuracy, F1 w~wariancie makro i~ważonym, strata logarytmiczna, metryki poszczególnych klas oraz wieloklasowe AUC one-vs-rest makro tam, gdzie jest określone. Główne etykiety decyzyjne powstają przez \(\arg\max\) prawdopodobieństw. Nie oznacza to optymalizacji kosztu aplikacyjnego. Moduł metryk zapisuje nieokreślone ROC AUC jako brak wartości, lecz nie zapisuje osobnego pola z~przyczyną. Dla precyzji, czułości i~F1 stosuje natomiast konwencję \texttt{zero\_division=0}: przy zerowym mianowniku zwraca zero. Rozróżnienie tych przypadków od zerowej wartości poprawnie określonej miary wymaga uwzględnienia liczności klas i~macierzy pomyłek.

W~badaniu rzadkich klas podstawową miarą jest average precision, dalej AP. W~standardowej definicji sumuje ona przyrosty czułości ważone odpowiadającą im precyzją; nie jest to trapezowe pole pod krzywą precision--recall. Uzupełnienie AUC o~tę perspektywę jest uzasadnione odmiennym zachowaniem ocen przy niezrównoważeniu~\cite{saito2015pr}. Interpretacja AP uwzględnia częstość klasy rzadkiej i~liczbę jej przykładów w~teście.

Rzadkie zdarzenie jest określone przez surową etykietę i~jej zapisane kodowanie. Zbiorcze binarne AP raportera dotyczy klasy zakodowanej jako~1, która nie musi oznaczać zdarzenia rzadkiego. W~badaniu~A kolumna prawdopodobieństw i~metryka odpowiadają semantycznie wskazanej klasie rzadkiej. Techniczne nazwy pól raportu i~ich mapowanie na populacje oceny opisano w~rozdziale~\ref{chap:srodowisko}.

\subsection{Badanie A: ważenie klas}
\label{sec:projekt-wazenie}

Badanie obejmuje Credit Card Fraud, Santander i~Amazon oraz metody MLP, TabM i~CatBoost. Dla każdego zestawu porównuje się uczenie nieważone i~ważone. Wagi klas są wyznaczane wyłącznie z~treningowej części danego foldu:
\begin{equation}
w_c=\frac{N_{\mathrm{tr}}}{K n_{c,\mathrm{tr}}},
\label{eq:wagi-klas}
\end{equation}
gdzie \(K\) oznacza liczbę klas, a~\(n_{c,\mathrm{tr}}\) --- liczbę treningowych przykładów klasy~\(c\). Zmianie podlega strata treningowa. Walidacja i~test zachowują oryginalny rozkład oraz nieważone metryki. Badanie nie obejmuje resamplingu ani odrębnego przeszukiwania hiperparametrów dla obu warunków.

Pozostałe hiperparametry są określane na podstawie wyboru OOF w~badaniu głównym. Porównanie wykorzystuje te same foldy i~trzy sparowane ziarna uczenia. W~ten sposób porównanie odpowiada na pytanie o~interwencję przy ustalonej procedurze, a~nie o~najlepszy wynik dwóch osobno strojonych rodzin. Układ obejmuje 54 oceny warunków i~432 dopasowania foldów: 288 neuronowych oraz 144 CatBoost na CPU.

Zgodnie z~analizą w~sekcji~\ref{sec:trudne-rezimy} wzrost czułości po ważeniu może wynikać z~przesunięcia punktu pracy bez poprawy rankingu. AP, ROC AUC i~nieważona strata logarytmiczna pozwalają rozdzielić te efekty. Precyzja, czułość, liczby FP i~FN oraz udział alarmów podlegają ocenie przy progach określonych na OOF. Wybór progu na walidacji definiuje badawczy punkt pracy; bez kosztów błędów nie stanowi optimum wdrożeniowego. Próg wybrany na OOF jest stosowany do średniej predykcji foldów, dlatego jego przenoszenie między tymi rozkładami ocen stanowi dodatkowe ograniczenie interpretacji.

\subsection{Badanie B: regularyzacja i liczność treningu}
\label{sec:projekt-regularyzacja}

Badanie obejmuje German Credit, Ionosphere i~Bioresponse. Punktem odniesienia jest procedura uczenia MLP oparta na opublikowanym rozwiązaniu i~ustalona niezależnie od wyników strojenia na tych trzech zadaniach. Procedura musi zawierać niezerowy dropout oraz niezerowy zanik wag. Nie należy utożsamiać jej automatycznie z~lokalnym bazowym MLP badania głównego. Przedmiotem porównania są pełna procedura, wariant bez dropoutu i~wariant bez zaniku wag. Brak wariantu usuwającego oba składniki ogranicza wnioski o~ich interakcji.

Dla każdego treningowego foldu tworzy się zagnieżdżone, stratyfikowane podzbiory o~udziałach 25\%, 50\% i~100\%. Redukcja obejmuje wyłącznie trening, z~ponownym dopasowaniem jego przekształceń. Walidacja i~test pozostają stałe. Definicja podzbioru obejmuje rzeczywiste liczności klas po zaokrągleniach; nominalny udział nie wystarcza do odtworzenia małych podzbiorów. Te same indeksy i~trzy sparowane ziarna są stosowane dla wszystkich trzech wariantów regularyzacji.

Projekt~\ref{fig:podzbiory} pozwala oddzielić interwencję regularyzacyjną od zmiany składu walidacji. Nie jest krzywą kosztu etykietowania: nawet najmniejszy trening korzysta z~niezmniejszonej walidacji. Układ obejmuje 81 ocen warunków i~648 neuronowych dopasowań foldów. Nie przeprowadza się odrębnego przeszukiwania 26 konfiguracji dla każdej komórki projektu.

\plannedfigure
{figures/ch04/zagniezdzone-podzbiory.pdf}
{Układ zagnieżdżonych podzbiorów treningowych w~badaniu regularyzacji; opracowanie własne.}
{fig:podzbiory}
{\textbf{Trening wewnętrznego foldu}\\
$\mathcal{T}_{25\%}\subset\mathcal{T}_{50\%}\subset\mathcal{T}_{100\%}$\\[0.8em]
Każdy podzbiór: pełna procedura MLP, bez dropoutu, bez zaniku wag.\\
Te same indeksy i~sparowane ziarna dla wariantów regularyzacji.\\[0.8em]
\textbf{Walidacja i~test}: stałe, poza redukcją treningu.\\
Przekształcenia cech dopasowywane na właściwym podzbiorze treningowym.}

Analiza generalizacji łączy testowe AUC z~historią uczenia oraz nieważoną stratą logarytmiczną na treningu i~walidacji. Luka jest definiowana dla tej samej sieci foldu, w~trybie ewaluacji z~wyłączonym dropoutem. Różnica między stratą treningową pojedynczej sieci a~testową stratą zespołu ośmiu sieci nie jest luką generalizacji tej samej procedury. Mniejsza luka sama w~sobie nie oznacza lepszej generalizacji; może towarzyszyć pogorszeniu obu ocen.

\subsection{Analiza porównawcza i niepewność}
\label{sec:analiza-plan}

Analiza rozdziela wielkość różnicy, jej zmienność oraz zakres wnioskowania. Podstawowy panel obejmuje T dla dziesięciu metod strojonych i~odrębnie oznaczony stały TabPFN. D i~T+E służą do analizy zmiany procedury. Każdy wynik pozostaje przypisany do jednego zadania, metryki i~warunku.

\begin{table}[tbp]
\centering\small
\caption{Powiązanie pytań badawczych z~porównaniami i~jednostkami analizy.}
\label{tab:pytania-analiza}
\begin{tabularx}{\linewidth}{P{0.09\linewidth}YY}
\toprule
Pytanie & Porównanie i~miara & Jednostka i~granica wniosku \\
\midrule
PB1 & Wyniki wszystkich metod; kontrasty między procedurami & Zbiór; celowo wybrane zadania \\
PB2 & D--T i~T--T+E; zmiana metryki oraz koszt & Ta sama para zbiór--metoda \\
PB3 & Ważenie względem kontroli; AP i~decyzje progowe & Sparowane ziarna w~trzech zadaniach \\
PB4 & Ablacja względem pełnej procedury; AUC, strata i~luka & Wspólny podzbiór i~ziarno w~trzech zadaniach \\
PB5 & Jakość oraz rozdzielone składniki kosztu & Zadanie, procedura i~określony sprzęt \\
\bottomrule
\end{tabularx}
\end{table}

Kontrasty o~bezpośredniej interpretacji metodologicznej to RealMLP\allowbreak{}--\allowbreak{}CatBoost, TabM\allowbreak{}--\allowbreak{}CatBoost oraz ModernNCA\allowbreak{}--\allowbreak{}KNN. Dla pary \(A,B\) różnica dodatnia oznacza korzyść~\(A\): \(\Delta_{\mathrm{AUC}}=\mathrm{AUC}_A-\mathrm{AUC}_B\), natomiast \(\Delta_{\log}=L_{\log,B}-L_{\log,A}\). Surowe różnice AUC i~straty nie podlegają uśrednianiu na wspólnej skali. Pomocnicze rangi można obliczać jedynie przy podaniu wspólnego zbioru kompletnych wyników i~reguły remisów.

Wnioskowanie o~przewadze między zadaniami wymaga określenia modelu losowego, rodziny porównań i~sposobu postępowania z~remisami. Test znaków dotyczy częstości kierunku różnicy, nie jej wielkości; równoczesne testowanie wielu par wymaga kontroli wielokrotności. Dobór celowy nie gwarantuje niezależności zadań ani modelu losowego uzasadniającego test. Jeżeli założeń nie da się obronić, analiza ma charakter opisowy. Niewielka liczba zadań ogranicza moc, a~brak istotności nie potwierdza równoważności.

Niepewność różnicy na pojedynczym zadaniu można analizować przez sparowany bootstrap zapisanych predykcji. Przy założeniu IID jednostką losowania jest wiersz, z~warstwowaniem po klasie. Dla Diabetes jednostkę stanowi pacjent ze wszystkimi jego pobytami. Te same repliki muszą służyć obu modelom. Liczba replik z~nieokreśloną metryką stanowi część opisu resamplingu. Takie przedziały są warunkowe względem ustalonego podziału i~wytrenowanych modeli; nie mierzą zmienności ponownego uczenia.

W~badaniach~A i~B przedmiotem analizy są wyniki wszystkich trzech sparowanych ziaren, średnia różnica i~jej zakres. Zakres nie jest przedziałem ufności. Foldy nie są traktowane jako niezależne replikacje, a~trzy zadania każdego badania nie uzasadniają przewagi populacyjnej. Pełne zestawienie jedenastu metod nie wymaga testowania wszystkich 55 par; niezgodne skale metryk pozostają rozdzielone.

Braki i~awarie wymagają jawnych mianowników. Analiza kompletnych par może pomijać najtrudniejsze zadania i~przez to zmieniać obraz porównania. Nieudanym uruchomieniom nie przypisuje się sztucznych metryk. Wielkość różnicy oraz jej znaczenie praktyczne pozostają odrębne od wartości~\(p\); bez kosztów aplikacyjnych nie ustala się wspólnej granicy użyteczności dla AUC, AP i~straty.

\subsection{Koszt obliczeń i odtwarzalność}
\label{sec:koszt-metodyka}

Nominalna macierz główna obejmuje \(12\cdot10\cdot26+12=3132\) zadań konfiguracji, 132 pary zbiór--metoda i~372 wiersze raportu dla przewidzianych warunków. Sześć strojonych procedur neuronowych wymaga nominalnie 14\,976 dopasowań foldów; stała inferencja TabPFN jest liczona osobno. Badania mechanizmów dodają 936 dopasowań neuronowych i~144 na CPU. Są to nominalne liczby dopasowań wynikające z~macierzy badania; rzeczywisty koszt zależy także od czasu pojedynczego dopasowania, powtórzeń i~niepowodzeń.

Koszt wyszukiwania i~suma czasów dopasowania konfiguracji wybranych do użycia opisują różne etapy procedury. Czas \texttt{search\_fit\_seconds} obejmuje udane dopasowania, natomiast \texttt{search\_wall\_seconds} uwzględnia zapisane próby, także zakończone błędem i~powtórzenia. Przerwanie bez znanego końca próby może uniemożliwić obliczenie łącznego czasu rzeczywistego. Brakująca wartość nie oznacza zerowego kosztu.

\texttt{selected\_inference\_seconds} sumuje zapisane przebiegi predykcji konfiguracji o~dodatniej wadze. Nie jest nowym pomiarem końcowej latencji całego zespołu i~nie obejmuje automatycznie wszystkich kosztów składania wyników. Czas wyboru z~pamięci podręcznej raportowany jest osobno. Koszt znalezienia kandydatów jest wspólny dla D, T i~T+E, dlatego sumowanie tych wierszy prowadziłoby do jego wielokrotnego naliczenia.

Opis warunków pomiaru obejmuje CPU, GPU, RAM, VRAM, liczbę wątków, rozmiar partii predykcji i~granice pomiarów. Ustawienia obejmują cztery wątki CPU. Limit partii predykcji~4096 jest przekazywany lokalnym adapterom i~wykorzystywany przez MLP, FT-Transformer oraz TabNet. Dla MLP i~FT-Transformer obowiązuje dodatkowo ograniczenie efektywnym rozmiarem partii uczenia. Natywne KNN, RealMLP, ModernNCA i~TabM oraz stały TabPFN nie otrzymują tego ustawienia; podział predykcji na partie określają te implementacje. W~przyjętym podziale zasobów sześć uczonych metod neuronowych korzysta z~GPU, a~odniesienia konwencjonalne z~CPU. Porównanie kosztu na takim podziale zasobów nie izoluje efektu architektury. Backend zapisuje pomiary pamięci w~bloku \texttt{memory\_usage} rekordu konfiguracji, a~nie w~kolumnach zbiorczego CSV. Śledzenie obejmuje wyłącznie dopasowanie: RSS procesu na CPU oraz pamięć przydzieloną i~rezerwowaną przez PyTorch na GPU. Interpretacja odczytu GPU zależy od znacznika dostępności śledzenia. Pomiary te nie określają pełnego zużycia VRAM procesu, szczytu całej próby ani pamięci testowej predykcji. Ta ostatnia wymaga osobnego pomiaru, szczególnie dla KNN i~TabPFN. Zakres pomiaru nie obejmuje zewnętrznego pretrenowania TabPFN, kosztu finansowego ani emisji.

W~FT-Transformerze efektywny rozmiar partii jest dodatkowo ograniczany na podstawie liczby cech i~głów uwagi. Dla $p$ predyktorów i~$h$ głów przyjmuje się $4h(p+1)^2$ bajtów na przykład dla pojedynczego tensora wyników uwagi w~formacie float32, z~uwzględnieniem tokenu klasyfikacyjnego. Z~przydziału 4~GiB wyznacza się limit liczby przykładów, zaokrąglany w~dół do potęgi dwójki, z~dolnym ograniczeniem tego limitu do ośmiu. Efektywna partia jest mniejszą z~wartości skonfigurowanej i~tak wyznaczonego limitu. Przy 1776 predyktorach limit wynosi 32 przykłady dla ośmiu głów oraz 64 dla czterech. Ograniczenie obowiązuje także podczas walidacji i~predykcji. Jest to heurystyka pamięciowa dotycząca jednego tensora, nie gwarancja zmieszczenia pełnego przebiegu w~VRAM. Ponieważ zmniejszenie partii zmienia również liczbę aktualizacji parametrów na epokę, efektywny rozmiar należy uwzględniać przy interpretacji zarówno kosztu, jak i~jakości modelu.

Tożsamość danych, zależności i~sprzętu wiąże pomiar z~konkretnym wykonaniem. Mechanizm zapisu tej informacji oraz kontroli zgodności opisano w~rozdziale~\ref{chap:srodowisko}.
